% ===== tr-02 块2：第 3 章 pp.15--18（原文页 15--18）=====
% [原文 p.15]
我们在此不给出这一紧化如何进行的具体细节；但只要把范围限制在拓扑（而非 $C^\infty$）流形范畴内，这并不特别困难。不过，考察我们需要在这一紧化中加入的奇异纤维是很有意思的。

若 $b\in\Gamma$ 不是 $\Gamma$ 的顶点，则 $f^{-1}(b)$ 同胚于 $I_1\times S^1$，其中 $I_1$ 表示 Kodaira $I_1$ 型椭圆曲线，即一个被捏住的环面（pinched torus）。

若 $b$ 是 $\Gamma$ 的顶点、且单值为第一类，则 $f^{-1}(b)=S^1\times S^1\times S^1/{\sim}$，其中 $(a,b,c)\sim(a',b',c')$ 当且仅当 $(a,b,c)=(a',b',c')$ 或 $a=a'=1$，这里 $S^1$ 等同于 $\C$ 中的单位圆。这是被捏住的环面在三维情形下的类比物，并且 $\chi(f^{-1}(b))=+1$。我们称这样的纤维为\emph{正纤维}。

若 $b$ 是 $\Gamma$ 的顶点、且单值为第二类，则 $f^{-1}(b)$ 可以描述为 $S^1\times S^1\times S^1/{\sim}$，其中 $(a,b,c)\sim(a',b',c')$ 当且仅当 $(a,b,c)=(a',b',c')$，或 $a=a'=1$ 且 $b=b'$，或 $a=a'$ 且 $b=b'=1$。该纤维的奇异轨迹是一个 8 字形（figure eight），并且 $\chi(f^{-1}(b))=-1$。我们称这样的纤维为\emph{负纤维}。

% [Tutor Note: 两处 $S^1\times S^1\times S^1/{\sim}$ 的等价关系不同，正是两类顶点的区别：第一类只允许把一个坐标（a）压成 1，第二类还允许 b 方向也退化；由此 Euler 特征数分别为 +1 与 −1。原文的 (a,b,c) 与 a',b',c' 上的撇号要逐一看清（500 dpi 页图复核）。]

于是我们看到 SYZ 对偶性的一个非常具体的局部推论：在紧化 $X(B)$ 与 $\check X(B)$ 中，正纤维与负纤维互换。当然，这导致如下观察：对 Calabi--Yau 三维流形，Euler 特征数在镜像对称下变号。

\begin{example}{3.4.5}
沿用例 3.4.3，[17] 中证明了 $\check X(B)$ 同胚于五次三维流形，而 $X(B)$ 同胚于镜像五次三维流形。
\end{example}

\begin{remark}{3.4.6}
Haase 与 Zharkov 在 [28] 中给出了同一个仿射结构的另一种描述。他们的构造的优点是显然对偶。换言之，在我们的构造中，我们可以互换 $\Delta$ 与 $\nabla$ 的角色，得到两个不同的仿射流形：$B_\Delta$ 是 $\partial\Delta^h$ 上的带奇点仿射结构，而 $B_\nabla$ 是 $\partial\nabla^{\check h}$ 上的带奇点仿射结构。这两个仿射流形是「对偶」的，这一点并不显然；至少就 $X(B_\nabla)$ 同胚于 $\check X(B_\Delta)$、$\check X(B_\nabla)$ 同胚于 $X(B_\Delta)$ 而言并不显然。在上面给出的构造中，这一点可由我们将在第 7 节讨论的离散 Legendre 变换推出。另一方面，我这里给出的构造将从本文后面讨论的退化构造中自然产生。
\end{remark}

% [Tutor Note: 注意原文这里两处上标的差别：$B_\Delta$ 对应 $\partial\Delta^h$（h 上不带 check），$B_\nabla$ 对应 $\partial\nabla^{\check h}$（h 上带 check）。照 500 dpi 页图复核，原文确实如此，故译文保留这一字形差异，不擅自统一。]

Ruan 在 [52] 中用辛流论证给出了环面簇中超曲面的\emph{拉格朗日环面纤维化}的一种描述，他的构造应当与辛流形 $\check X(B_0)$ 的一个\emph{辛}紧化相吻合。在三维情形下，这样的辛紧化目前已经由 Ricardo Castaño-Bernard 与 Diego Matessi [7] 构造出来。若把这一紧化应用于这里描述的带奇点仿射流形，所得的辛流形应当与相应的环面超曲面辛同胚，但这一点尚未得到证明。

我还要指出，上述三维情形中的显式紧化可以在所有维数中进行，并且将在 [21] 中这样做。我们将在那里一个远为一般的框架中证明：这些紧化随后同胚于所期待的 Calabi--Yau 流形。

\section{SYZ 猜想的问题以及如何绕过它们}

上一节表明，SYZ 猜想在纯拓扑层面上为镜像对称给出了一个漂亮的描述。这本身往往是有用的，但不幸的是，它还不足以触及镜像对称中真正有趣的那些方面，例如瞬子修正。不过，有一段时间，我们中的许多人曾希望我们刚刚看到的强版本对偶性在特殊拉格朗日层面上也成立。这将意味着：一对镜像流形 $X,\check X$ 会拥有特殊拉格朗日环面纤维化 $f:X\to B$ 与 $\check f:\check X\to B$，其判别轨迹余维为 2，且 $f$ 与 $\check f$ 的判别轨迹重合。于是这两个纤维化在判别轨迹之外是对偶的。

% [原文 p.16]
在非紧环面簇 $X$ 上，确实存在具有这种行为的特殊拉格朗日纤维化的例子。特别地，若 $\dim X=n$，$X$ 上有一个保持全纯 $n$-形式的 $T^{n-1}$ 作用，并且 $X$ 还带有一个在该作用下不变的 Ricci 平坦度量，则 $X$ 会有一个非常漂亮的、判别轨迹余维为 2 的特殊拉格朗日纤维化。（见 [18] 与 [14]。）然而，Dominic Joyce（[33] 及其所引的其他论文）开始研究一些三维的 $S^1$-不变例子，并发现了相当不同的行为。有一个论证表明：若特殊拉格朗日纤维化是 $C^\infty$ 的，则判别轨迹将是（Hausdorff 意义下）余维 2 的。然而，Joyce 发现的例子不可微，而只是分段可微的，并且其判别轨迹是余维 1 的：

\begin{example}{3.5.1}
定义 $F:\C^3\to\R\times\C$ 为 $F(z_1,z_2,z_3)=(a,c)$，其中 $2a=|z_1|^2-|z_2|^2$，且
\[
c=\begin{cases}
z_3 & a=z_1=z_2=0\\
z_3-\bar z_1\bar z_2/|z_1| & a\ge0,\ z_1\ne0\\
z_3-\bar z_1\bar z_2/|z_2| & a<0.
\end{cases}
\]
\end{example}

容易看出：若 $a\ne0$，则 $F^{-1}(a,c)$ 同胚于 $\R^2\times S^1$；而若 $a=0$，则 $F^{-1}(a,c)$ 是 $T^2$ 上的一个锥：本质上，$\R^2\times S^1$ 中有一份 $S^1$ 塌缩成一个点。此外，该映射的所有纤维都是特殊拉格朗日的，而它显然只是分段光滑的。其判别轨迹就是由 $a=0$ 给出的整个平面。

% [Tutor Note: 例 5.1 的三段定义中，第一段的条件是 a=z_1=z_2=0（不是单独的 a=0），第二段要求 a≥0 且 z_1≠0，第三段要求 a<0；a=0 但 (z_1,z_2)≠(0,0) 的情形落在第二段（若 z_1≠0）。三段合起来才覆盖 ℂ³。]

这个例子迫使人们重新评估 SYZ 猜想的强形式。在后续工作中，Joyce 找到了证据，支持三维一般特殊拉格朗日纤维化的一个更可能的图景。判别轨迹不再是余维 2 的图，而将是一个余维 1 的「团块」（blob）。典型地，奇异纤维的奇点之并会是一张 Riemann 面，而它将映入 $B$ 中一个变形虫状（amoeba）的集合，即判别轨迹看起来像下图中的右边而不是左边，并且将是旧的余维 2 判别轨迹图景的一个「加厚」。

% [原文 p.17]
\figimg[0.62]{p3_p17.png}

% [Tutor Note: 原文 p.17 页首此图无图号、无图注；p.16 末句 "looks like the picture on the right rather than the left" 即指此图（左=细的余维 2 判别轨迹，右=加厚的余维 1 判别轨迹）。译文按原图所在页插入，不加图号。]

% [Tutor Note: 图后正文说「the two basic types of vertices we saw in the previous section」：指 p.14 的两类三价顶点（二维面内部型与 ∇^{h} 棱上型），在加厚图景中这两类顶点附近的判别轨迹形态不同。]

Joyce 作了一些进一步的论证，暗示这个加厚的判别轨迹在上一节所见的两类基本顶点的邻域中必定看起来有根本的不同，而这两类顶点正如图片中所描绘的那样出现。于是，上面提到的强形式对偶性------即我们期待镜像对上的特殊拉格朗日纤维化的判别轨迹相同------不能成立。若情况如此，人们就需要用更弱的形式来取代这一强形式对偶性。

重新表述 SYZ 猜想的最佳方式很可能是一个极限形式。我们目前所理解的镜像对称与 Calabi--Yau 流形的退化有关。给定圆盘 $D$ 上的一族平坦族 $f:\mathcal{X}\to D$，其中 $0$ 上的纤维 $\mathcal{X}_0$ 是奇异的，而所有其他纤维都是 $n$ 维 Calabi--Yau 流形；若单值变换 $T:H^n(\mathcal{X}_t,\Q)\to H^n(\mathcal{X}_t,\Q)$（$t\in D$ 非零）满足 $(T-I)^{n+1}=0$ 但 $(T-I)^n\ne0$，则称该族是\emph{极大幂零}的。镜像应当与 Calabi--Yau 流形的极大幂零退化相关联，这是镜像对称的一个标准事实。特别地，给定某个 Calabi--Yau 流形的单个复模空间中的两个不同的极大幂零退化，人们可能得到不同的镜像流形。有时这些不同的镜像流形是双有理等价的（如 [1] 所研究），有时则确实是不同的，见 [51]。

% [Tutor Note: 极大幂零退化的定义里 $(T-I)^{n+1}=0$ 与 $(T-I)^n\ne0$ 要一起读：幂零指数恰为 n+1（这正是「极大」的含义），等价的说法是单值算子 T 在 H^n 上有唯一的 Jordan 块、其特征值全为 1。]

% [Tutor Note: 原文写 "$t\in D$ non-zero"，严格说是 $t\in D\setminus\{0\}$（圆盘去掉原点后单值才有定义）。译文照原文措辞，此处附带说明。]

我们回顾 Gromov--Hausdorff 收敛的定义，即度量空间序列收敛的一个概念。

\begin{definition}{3.5.2}
设 $(X,d_X)$、$(Y,d_Y)$ 是两个紧度量空间。若存在映射 $f:X\to Y$ 与 $g:Y\to X$（不必连续），使得对所有的 $x_1,x_2\in X$ 有
\[
|d_X(x_1,x_2)-d_Y(f(x_1),f(x_2))|<\epsilon,
\]
并且对所有的 $x\in X$ 有
\[
d_X(x,g\circ f(x))<\epsilon,
\]
且关于 $Y$ 的两个对称性质也成立，则称 $X$ 与 $Y$ 之间的 Gromov--Hausdorff 距离至多为 $\epsilon$。Gromov--Hausdorff 距离 $d_{GH}(X,Y)$ 是所有这些 $\epsilon$ 的下确界。
\end{definition}

% [Tutor Note: 定义 5.2 只显式写出了一侧的两个条件（f 近似等距、g∘f 近似恒等），原文用 "the two symmetric properties for Y hold" 一句把另一侧（g 也近似等距、f∘g 近似恒等）略去，译文照原文保留这一省略。]

由 Gromov 的结果（例如见 [49]，第 281 页，推论 1.11）可知，直径 $\le C$ 的紧 Ricci 平坦流形之空间是预紧的
% [原文 p.18]
（关于 Gromov--Hausdorff 距离），即这样的流形的任何序列都有一个子序列按 Gromov--Hausdorff 距离收敛到一个度量空间。这个度量空间可能相当糟糕；这完全在代数几何的疆域之外！然而，这引出了下面这个自然的问题。给定 Calabi--Yau 流形的一个极大幂零退化 $\mathcal{X}\to D$，取收敛到 0 的序列 $t_i\in D$，并考虑序列 $(\mathcal{X}_{t_i},g_{t_i})$，其中 $g_{t_i}$ 是选取的 Ricci 平坦度量，使得 $\operatorname{Diam}(g_{t_i})$ 保持有界。$(\mathcal{X}_{t_i},g_{t_i})$ 的 Gromov--Hausdorff 极限是什么，或者某个收敛子序列的极限是什么？

\begin{example}{3.5.3}
考虑一族退化的椭圆曲线，其周期为 1 与 $\frac{1}{2\pi i}\log t$。若让 $t$ 沿正实轴趋于 0，则我们可以直接把它看作一族椭圆曲线 $\mathcal{X}_\alpha$，其周期为 1 与 $i\alpha$，其中 $\alpha\to\infty$。若取 $\mathcal{X}_\alpha$ 上的标准欧氏度量 $g$，则 $\mathcal{X}_\alpha$ 的直径无界。为得到有界直径，我们把 $g$ 换成 $g/\alpha^2$；等价地，我们可以保持 $\C$ 上的 $g$ 不变，而把椭圆曲线的周期改为 $1/\alpha$ 与 $i$。于是可以看清楚：这样一列椭圆曲线的 Gromov--Hausdorff 极限是一个圆 $\R/\Z$。
\end{example}

% [Tutor Note: 原文把周期写成 $1/2\pi i\log t$ 的分式形式，即 $\frac{1}{2\pi i}\log t$（分母是 2πi），而不是 $\frac{1}{2\pi i\log t}$；文字层抽成 "2πi log t" 时分母归属极易抄错，须对照 p.18 页图。]

这个简单的例子引出了关于极大幂零退化的第一个猜想，它分别由我本人与 Wilson（[26]）以及 Kontsevich 与 Soibelman（[38]）独立提出。

\begin{conjecture}{3.5.4}
设 $\mathcal{X}\to D$ 是单连通、具有全 $SU(n)$ 和乐的 Calabi--Yau 流形的极大幂零退化，$t_i\in D$ 且 $t_i\to0$；设 $g_i$ 是 $\mathcal{X}_{t_i}$ 上规范化为固定直径 $C$ 的 Ricci 平坦度量。则 $(\mathcal{X}_{t_i},g_i)$ 的某个收敛子序列收敛到一个度量空间 $(X_\infty,d_\infty)$，其中 $X_\infty$ 同胚于 $S^n$。此外，$d_\infty$ 由 $X_\infty\setminus\Gamma$ 上的一个 Riemann 度量诱导，其中 $\Gamma\subseteq X_\infty$ 是一个余维 2 的集合。
\end{conjecture}

这里极限的拓扑取决于非奇异纤维 $\mathcal{X}_t$ 的性质；例如，若 $\mathcal{X}_t$ 改为超 Kähler 的，则我们会期待极限是一个射影空间。另外，即使在具有全 $SU(n)$ 和乐的情形下，若 $\mathcal{X}_t$ 不是单连通的，我们也会期待出现诸如 $\Q$-同调球之类的极限。

猜想 3.5.4 直接受 SYZ 猜想的启发。假设我们已有特殊拉格朗日纤维化 $f_i:\mathcal{X}_{t_i}\to B_i$。那么当趋近极大幂零退化时，可以看出这些纤维化的纤维体积趋于零。这提示这些纤维会坍缩，从而有望把底留作极限。

这个猜想由我本人与 Wilson 对 K3 曲面在 [26] 中证明。证明依赖于对 Ricci 平坦度量在极限中行为的细致分析，也依赖于特殊拉格朗日纤维化的奇异纤维附近 Ricci 平坦度量的显式局部模型的存在性。

由 SYZ 给出的这一猜想的动机，还提供了猜想的一个极限形式。要证明 Calabi--Yau 流形上特殊拉格朗日纤维化的存在性，问题数不胜数。甚至极大幂零退化附近单个特殊拉格朗日环面的存在性也未知，但我们期待当我们趋近极大幂零点时会更容易找到它们。此外，即使我们找到了一个特殊拉格朗日环面，我们知道它在一个 $n$ 维族中移动，但我们并不知道它的形变是否会填满整个流形。另外，即使填满，也不能保证我们得到流形的一个叶状结构：相邻的特殊拉格朗日子流形可能相交。
